\section{The secant family and the identity}\label{lv:sec:identity}

On the whole line no symmetric scale family can be complete, because odd statistics have mean zero, so the dilates are folded to $(0,\infty)$. The symmetric Meixner density of order $t>0$ is proportional to $|\Gamma(\frac t2+\frac{ix}2)|^2$ (Section~\ref{rz:sec:meixner}), and $t=1$ gives the hyperbolic secant density because $|\Gamma(\frac12+iy)|^2=\pi/\cosh(\pi y)$. Folded to $x>0$ it is $p_\lambda(x)=\frac{2\lambda}\pi\sech(\lambda x)$, and a family $\{p_\lambda:\lambda\in\Lambda\}$ is complete when $\E_\lambda g(X_\lambda)=0$ for all $\lambda\in\Lambda$ forces $g=0$ almost everywhere. The smallest natural parameter set is $\Lambda=\mathbb N$, so completeness says that the dilates $\sech(nx)$ separate functions; this is the integer case of the question of Lai and Vere-Jones \cite{LVJ} as recorded in Chapter~\ref{chap:rzeta}.

\begin{lemma}\label{lv:lem:exp}
For $y>0$, $\sech y=2\sum_{m\ge1}\chi_4(m)e^{-my}$, and $\sum_{n\ge1}\sech(ny)\le2/(e^y-1)$.
\end{lemma}

\begin{proof}
The first identity is the geometric series of Proposition~\ref{rz:prop:lvj}. For the bound use $\sech(ny)\le2e^{-ny}$.
\end{proof}

\begin{lemma}\label{lv:lem:mellin}
For $\Re s>0$ and $\lambda>0$, $\int_0^\infty x^{s-1}\sech(\lambda x)\,dx=2\lambda^{-s}\Gamma(s)\beta(s)$.
\end{lemma}

\begin{proof}
For $\lambda=1$ integrate Lemma~\ref{lv:lem:exp} termwise for $\Re s>1$, where $\sum_m\int_0^\infty x^{\sigma-1}e^{-mx}dx=\Gamma(\sigma)\zeta(\sigma)<\infty$, and continue analytically; the integral is holomorphic for $\Re s>0$ and $\beta$ is entire. The substitution $y=\lambda x$ gives general $\lambda$.
\end{proof}

By Proposition~\ref{rz:prop:lvj}, $e^{-x}=\frac12\sum_{n\ge1}\mu(n)\chi_4(n)\sech(nx)$ with absolute convergence for $x>0$, and replacing $x$ by $kx$,
\begin{equation}\label{lv:eq:LVJk}
e^{-kx}=\frac12\sum_{n\ge1}\mu(n)\chi_4(n)\sech(nkx),\qquad k\in\mathbb N,
\end{equation}
and every function on the right is again one of the dilates. In Mellin language the identity is the factorisation $\Gamma(s)=\frac12\cdot2\Gamma(s)\beta(s)\cdot\beta(s)^{-1}$ with $\beta(s)^{-1}=\sum\mu(n)\chi_4(n)n^{-s}$.

\begin{theorem}[Lerch]\label{lv:thm:lerch}
Let $\int_0^\infty|h(x)|e^{-x}dx<\infty$. If $\int_0^\infty e^{-kx}h(x)\,dx=0$ for all $k\in\mathbb N$, then $h=0$ almost everywhere.
\end{theorem}

\begin{proof}
With $u=e^{-x}$ and $q(u)=h(-\log u)$ we have $q\in L^1(0,1)$ and $\int_0^1u^jq(u)\,du=0$ for $j\ge0$. The measure $q\,du$ annihilates the polynomials, hence $C[0,1]$ by Weierstrass, hence $q=0$.
\end{proof}

\begin{theorem}[completeness in $L^1$]\label{lv:thm:L1}
Let $g$ satisfy
\begin{equation}\label{lv:eq:L1cond}
\int_0^1\frac{|g(x)|}x\,dx+\int_1^\infty|g(x)|e^{-x}dx<\infty .
\end{equation}
If $\int_0^\infty g(x)\sech(nx)\,dx=0$ for all $n\in\mathbb N$, then $g=0$ almost everywhere. The same holds with $\mathbb N$ replaced by a set $\Lambda\subset(0,\infty)$ containing an interval, provided the integrals $\int g\sech(\lambda x)\,dx$, $\lambda\in\Lambda$, converge absolutely and vanish.
\end{theorem}

\begin{proof}
By Lemma~\ref{lv:lem:exp}, $\sum_n|\sech(nkx)|\le2/(e^{kx}-1)$, which is $O(1/x)$ at the origin and $O(e^{-x})$ at infinity, so \eqref{lv:eq:L1cond} justifies integrating \eqref{lv:eq:LVJk} termwise against $g$. Every term vanishes, hence $\int ge^{-kx}dx=0$ for all $k$, and Theorem~\ref{lv:thm:lerch} applies. For the interval clause, let $I\subset\Lambda$ be an interval and $\lambda_0$ a point of its interior. Absolute convergence at $\lambda_0$ gives $\int_1^\infty|g|e^{-\lambda_0x}dx<\infty$, because $\sech(\lambda_0x)\ge e^{-\lambda_0x}$. The function $F(z)=\int_0^\infty g(x)\sech(zx)\,dx$ is then holomorphic on $\Re z>\lambda_0$ (on compact subsets the large $x$ part is dominated by $C|g(x)|e^{-(\lambda_0+\delta)x}$ and the small $x$ part by $C|g(x)|$), and it vanishes on a subinterval of $I$ to the right of $\lambda_0$, hence on the half plane. Take an integer $K>\max(1,\lambda_0)$. The integrals $\int g(x)\sech(nkx)\,dx$ vanish for all integers $n\ge1$ and $k\ge K$, so \eqref{lv:eq:LVJk} (again integrated termwise, the majorant being $|g|\,2/(e^{kx}-1)$ with $k\ge K$) gives $\int g(x)e^{-kx}dx=0$ for $k\ge K$. Apply Theorem~\ref{lv:thm:lerch} to $h(x)=g(x)e^{-(K-1)x}$, whose moments of orders $1,2,\dots$ are these integrals.
\end{proof}

The proof of the first part uses only $\sum_{d\mid N}\mu(d)=[N=1]$, the complete multiplicativity of $\chi_4$ and $|\mu|\le1$. The condition $\int_0^1|g|/x<\infty$ is the $L^1$ shadow of the half plane $\Re s>1$ in which the M\"obius series converges absolutely. The theorem is proved here in the form stated; it is not claimed to be the form in which the question was posed in \cite{LVJ}.

\section{Weighted spaces and the Mellin line}\label{lv:sec:weights}

Let $\mathcal W_c=L^2((0,\infty),x^{2c-1}dx)$. The change of variable $x=e^u$ and Plancherel give
\begin{equation}\label{lv:eq:plancherel}
\begin{gathered}\langle f,g\rangle_c=\int_0^\infty f\bar g\,x^{2c-1}dx=\frac1{2\pi}\int_{-\infty}^\infty F(c+it)\overline{G(c+it)}\,dt,\\ F(s)=\int_0^\infty f(x)x^{s-1}dx .
\end{gathered}\end{equation}
So $\mathcal W_c$ is the space of functions whose Mellin transform lives on the line $\Re s=c$. For $c>0$ both $e^{-x}$ and every $\sech(\lambda x)$ lie in $\mathcal W_c$, and dilation acts by $\|f(k\,\cdot)\|_c=k^{-c}\|f\|_c$.

\begin{lemma}[reduction]\label{lv:lem:reduction}
For $c>0$, or for $\mathcal W_a\cap\mathcal W_b$ with $0<a<b$ and norm $\|f\|_a+\|f\|_b$, the exponentials $\{e^{-kx}\}_{k\in\mathbb N}$ are complete. Hence a system whose linear span is invariant under the dilations $f\mapsto f(k\,\cdot)$, $k\in\mathbb N$, is complete as soon as its closed span contains $e^{-x}$.
\end{lemma}

\begin{proof}
A functional on $\mathcal W_c$ is $f\mapsto\int fh\,dx$ with $h=\bar gx^{2c-1}$, $g\in\mathcal W_c$, and $\int|h|e^{-x}\le\|g\|_c\|e^{-x}\|_c$; apply Theorem~\ref{lv:thm:lerch}. On the intersection the dual is the sum of the two duals (Hahn Banach) and the same estimate applies to each piece. If $f_j\to e^{-x}$ with $f_j$ in the span, then $f_j(k\,\cdot)\to e^{-k\,\cdot}$ in the same space and $f_j(k\,\cdot)$ is again in the span.
\end{proof}

Write $S_N(x)=\frac12\sum_{n\le N}\mu(n)\chi_4(n)\sech(nx)$. By Lemma~\ref{lv:lem:mellin} and \eqref{lv:eq:plancherel},
\begin{equation}\label{lv:eq:error}
\|S_N-e^{-x}\|_c^2=\frac1{2\pi}\int_{-\infty}^\infty|\Gamma(c+it)|^2\,\bigl|\beta(c+it)\mathcal A_N(c+it)-1\bigr|^2\,dt,
\end{equation}
so convergence of the M\"obius secant series in $\mathcal W_c$ is the statement that $\mathcal A_N\to1/\beta$ on the line $\Re s=c$ in a norm weighted by $|\Gamma(c+it)|^2\asymp|t|^{2c-1}e^{-\pi|t|}$.

\begin{theorem}[the classical range]\label{lv:thm:c1}
For $c\ge1$, $S_N\to e^{-x}$ in $\mathcal W_c$.
\end{theorem}

\begin{proof}
For $c>1$ use $|\mathcal A_N|\le\zeta(c)$, $|\beta|\le\zeta(c)$ and dominated convergence in \eqref{lv:eq:error}. For $c=1$ the prime number theorem for $\mu\chi_4$ in the form $M_\chi(x)\ll x\exp(-a\sqrt{\log x})$ and partial summation,
\begin{equation}\label{lv:eq:ps}
\mathcal A_N(s)=\frac{M_\chi(N)}{N^s}+s\int_1^N\frac{M_\chi(u)}{u^{s+1}}du,
\end{equation}
give $|\mathcal A_N(1+it)-1/\beta(1+it)|\le C(1+|t|)\varepsilon_N$ with $\varepsilon_N\to0$ (the integral converges on $\Re s=1$ by the same bound for $M_\chi$); with $|\beta(1+it)|\ll\log(2+|t|)$, \eqref{lv:eq:error} tends to zero.
\end{proof}

\section{Completeness in one weight is unconditional}\label{lv:sec:single}

The M\"obius series is one particular approximation of $e^{-x}$. Completeness asks only that some approximation exists, and in a single weight this needs no arithmetic. The tool is Carleman's formula for a half plane (Titchmarsh \cite{Titchmarsh39}; Boas \cite{Boas54}): if $f$ is analytic in a neighbourhood of $\Re w\ge0$, $f\not\equiv0$, with zeros $r_ke^{i\theta_k}$, then for $R>1$
\begin{multline}\label{lv:eq:carleman}
\sum_{1<r_k<R}\Bigl(\frac1{r_k}-\frac{r_k}{R^2}\Bigr)\cos\theta_k=\frac1{\pi R}\int_{-\pi/2}^{\pi/2}\log|f(Re^{i\theta})|\cos\theta\,d\theta\\
+\frac1{2\pi}\int_1^R\Bigl(\frac1{y^2}-\frac1{R^2}\Bigr)\log|f(iy)f(-iy)|\,dy+O_f(1).
\end{multline}
The form of the formula was checked against the standard statement; the section numbers of the two books are not claimed.
% UNCHECKED: Carleman formula section numbers (Titchmarsh 1939 Section 3.7, Boas 1954 Section 1.2) are the source's.

\begin{lemma}[uniqueness on a divergent set]\label{lv:lem:uniq}
Let $\Phi$ be analytic on $\Re w>0$ with $|\Phi(w)|\le C(|w|/\Re w)^B$ for some $B\ge0$. If $\Phi(\lambda)=0$ for every $\lambda$ in a set $\Lambda\subset[1,\infty)$ with $\sum_{\lambda\in\Lambda}1/\lambda=\infty$, then $\Phi\equiv0$.
\end{lemma}

\begin{proof}
If $\Lambda$ has a finite accumulation point, $\Phi\equiv0$ by the identity theorem, so let the points of $\Lambda$ be isolated. Fix $\varepsilon\in(0,1)$ and apply \eqref{lv:eq:carleman} to $f(w)=\Phi(w+\varepsilon)$, which is analytic on $\Re w>-\varepsilon$, assuming $f\not\equiv0$. Its zeros include $\lambda-\varepsilon$, $\lambda\in\Lambda$, on the positive axis (the finitely many with $\lambda-\varepsilon\le1$ lie outside the range of the sum and cost $O(1)$). The other zeros of $f$ contribute nonnegative terms ($\cos\theta_k\ge0$, $r_k<R$). Since $\frac1r-\frac r{R^2}\ge\frac34\cdot\frac1r$ for $r<R/2$, the left side is at least $\frac34\sum_{\lambda<R/2}1/\lambda-O(1)\to\infty$. On the right, for $w=Re^{i\theta}$ we have $|w+\varepsilon|/\Re(w+\varepsilon)\le2/\cos\theta$, so $\log|f|\le C'+B\log(1/\cos\theta)$ and the first integral is at most $O(1/R)$, because $\int\log(1/\cos\theta)\cos\theta\,d\theta<\infty$. On the imaginary axis $|f(\pm iy)|\le C(1+|y|/\varepsilon)^B$, so the second integral is at most $C_\varepsilon\int_1^\infty\log(2+y)\,y^{-2}dy<\infty$, the weight $\frac1{y^2}-\frac1{R^2}$ being nonnegative. The right side is bounded above, a contradiction. Hence $f\equiv0$.
\end{proof}

\begin{theorem}[completeness in one weight]\label{lv:thm:single}
Let $c>0$ and let $\Lambda\subset[1,\infty)$ satisfy $\sum_{\lambda\in\Lambda}1/\lambda=\infty$. Then $\{\sech(\lambda x)\}_{\lambda\in\Lambda}$ is complete in $\mathcal W_c$. In particular the integer dilates are complete in every $\mathcal W_c$, and so are the prime dilates $\{\sech(px)\}_{p\text{ prime}}$.
\end{theorem}

\begin{proof}
Let $g\in\mathcal W_c$ be orthogonal to every $\sech(\lambda x)$, and let $G$ be its Mellin transform on $\Re s=c$. By Lemma~\ref{lv:lem:mellin} and \eqref{lv:eq:plancherel},
\[
\int_{-\infty}^\infty H(t)\lambda^{-it}dt=0\quad(\lambda\in\Lambda),\qquad H(t)=\Gamma(c+it)\beta(c+it)\overline{G(c+it)} .
\]
Stirling and the convexity bound give $|\Gamma(c+it)\beta(c+it)|\le C(1+|t|)^Ae^{-\pi|t|/2}$, and $G\in L^2$, so $H\in L^1$. For $w=|w|e^{i\varphi}$ with $|\varphi|<\pi/2$ define
\[
\Phi(w)=\int_{-\infty}^\infty H(t)\,w^{-it}\,dt,\qquad|w^{-it}|=e^{t\varphi}.
\]
By Cauchy Schwarz,
\[
|\Phi(w)|\le C\|G\|_2\Bigl(\int(1+|t|)^{2A}e^{-(\pi-2|\varphi|)|t|}dt\Bigr)^{1/2}\le C'\Bigl(\frac\pi2-|\varphi|\Bigr)^{-A-1/2},
\]
and since $\frac\pi2-|\varphi|\ge\cos\varphi=\Re w/|w|$, $\Phi$ satisfies the growth bound of Lemma~\ref{lv:lem:uniq} with $B=A+\frac12$. It is analytic and vanishes on $\Lambda$, so $\Phi\equiv0$. On the positive axis $\Phi(e^\xi)=\int H(t)e^{-i\xi t}dt$ is the Fourier transform of $H\in L^1$, so $H=0$ almost everywhere. The zeros of $\beta$ on the line $\Re s=c$ form a discrete set, hence $G=0$ almost everywhere and $g=0$.
\end{proof}

\begin{remark}[the threshold, heuristic]\label{lv:rem:threshold}
If the Mellin transform of a kernel is bounded by $e^{-\alpha|t|}$ times a power, the function $\Phi$ above is analytic at least in the sector $|\arg w|<\alpha$. The map $w\mapsto w^{\pi/(2\alpha)}$ sends the sector to a half plane and $\lambda$ to $\lambda^{\pi/(2\alpha)}$, so the Carleman sum becomes $\sum\lambda^{-\pi/(2\alpha)}$; for $\Lambda=\mathbb N$ this diverges if and only if $\alpha\ge\pi/2$. The gamma factor of $\sech$ decays at precisely $\alpha=\pi/2$. This is the threshold of the method: failure of completeness for $\alpha<\pi/2$ is not shown, and the sharpness question is Open Problem~\ref{lv:prob:sharp}.
\end{remark}

\begin{remark}[M\"untz and Sz\'asz]
For the kernel $e^{-x}$ itself ($\beta\equiv1$) the proof of Theorem~\ref{lv:thm:single} gives the sufficiency half of the M\"untz and Sz\'asz theorem: under $u=e^{-x}$, completeness of $\{e^{-\lambda x}\}$ in $\mathcal W_{1/2}$ is completeness of $\{u^{\lambda-1/2}\}$ in $L^2(0,1)$, and divergence of $\sum1/\lambda$ is the classical sufficient condition. Lerch's theorem is the case $\Lambda=\mathbb N$.
\end{remark}

Three consequences are worth stating plainly. First, taking $c=\frac12$, the integer dilates of $\sech$ are complete in $L^2(0,\infty)$, even though $\beta$ has infinitely many zeros on the very line $\Re s=\frac12$ on which the Mellin transform lives. Second, no property of $\mu$ is needed for completeness in one weight. Third, the same proof applies to any kernel whose Mellin transform on the line is $\Gamma(s)$ times a polynomially bounded function vanishing only on a null set of the line, in particular to the zeta kernel of Section~\ref{lv:sec:zeta}. Theorem~\ref{lv:thm:single} therefore settles the one weight form of the completeness question recorded in Chapter~\ref{chap:rzeta}, unconditionally; only the two weight form is equivalent to a zero free half plane (Theorem~\ref{lv:thm:two}).

\section{The M\"obius route in the strip $\frac12<c<1$}\label{lv:sec:strip}

Where arithmetic does enter is in the explicit approximation by the M\"obius series.

\begin{lemma}[partial summation]\label{lv:lem:ps}
Suppose $M_\chi(x)=O_\varepsilon(x^{\theta+\varepsilon})$ for every $\varepsilon>0$, with $\frac12\le\theta<1$. Then for $\sigma=\Re s>\theta$, $t=\Im s$ and $0<\varepsilon<\sigma-\theta$, with constants depending on $\sigma$ and $\varepsilon$,
\[
|\mathcal A_N(s)|\le C(1+|t|),\qquad|\mathcal A_N(s)-1/\beta(s)|\le C(1+|t|)N^{\theta+\varepsilon-\sigma}.
\]
\end{lemma}

\begin{proof}
In \eqref{lv:eq:ps} the integral converges absolutely for $\sigma>\theta$ to an analytic function equal to $1/(s\beta(s))$ for $\sigma>1$, hence throughout; the boundary term and the tail $|s|\int_N^\infty u^{\theta+\varepsilon-\sigma-1}du$ give the rate.
\end{proof}

\begin{theorem}[the strip]\label{lv:thm:strip}
Under the hypothesis of Lemma~\ref{lv:lem:ps}, for every $c\in(\theta,1)$,
\[
\|S_N-e^{-x}\|_c=O_\varepsilon(N^{\theta-c+\varepsilon}).
\]
Under the generalised Riemann hypothesis for $\beta$ this holds for every $c>\frac12$.
\end{theorem}

\begin{proof}
Insert Lemma~\ref{lv:lem:ps} in \eqref{lv:eq:error}. The integral $\int|\Gamma\beta|^2(1+|t|)^2dt$ on the line $\Re s=c$ is finite by Stirling and convexity.
\end{proof}

\begin{proposition}[classical equivalences]\label{lv:prop:classical}
For $\frac12\le\theta<1$ the following are equivalent: (a) $M_\chi(x)=O(x^{\theta+\varepsilon})$ for all $\varepsilon>0$; (b) $\sum\mu(n)\chi_4(n)n^{-s}$ converges for $\Re s>\theta$; (c) $\beta(s)\ne0$ for $\Re s>\theta$.
\end{proposition}

\begin{proof}
(a) implies (b) by partial summation. (b) implies (c): the sum is analytic on $\Re s>\theta$ and equals $1/\beta$ for $\Re s>1$, so $1/\beta$ continues analytically there and $\beta$ has no zeros. (c) implies (a) is Littlewood's theorem for $L(s,\chi_4)$: a zero free half plane gives $1/\beta(\sigma+it)\ll|t|^{\varepsilon}$ for $\sigma>\theta$ by the Borel and Carath\'eodory argument, and Perron's formula then bounds $M_\chi$; compare Titchmarsh \cite[Section 14.25]{Titchmarsh} for $\zeta$, whose proof carries over to the character.
\end{proof}
% UNCHECKED: the general theta case of (c) implies (a) is quoted from the zeta case (Titchmarsh Section 14.25, Montgomery and Vaughan Chapter 13); the sections were not read for this book.

\section{Two weights: where the zeros are felt}\label{lv:sec:two}

The one weight proof breaks down for an intersection of two weights for a structural reason. If $f\in\mathcal W_a\cap\mathcal W_b$ and $a<\sigma<b$, splitting at $x=1$ and using Cauchy Schwarz with weight $a$ on $(0,1)$ and weight $b$ on $(1,\infty)$ gives
\begin{equation}\label{lv:eq:pointeval}
\Bigl|\int_0^\infty f(x)x^{s-1}dx\Bigr|\le\frac{\|f\|_a}{\sqrt{2(\sigma-a)}}+\frac{\|f\|_b}{\sqrt{2(b-\sigma)}} .
\end{equation}
So evaluation of the Mellin transform at any point of the open strip is continuous, and zeros of $\beta$ in the strip are inherited by limits. In a single weight there is no strip and evaluation off the line is not continuous.

\begin{proposition}[necessity]\label{lv:prop:nec}
Let $0<a<b$. If $e^{-x}$ is a limit in $\mathcal W_a\cap\mathcal W_b$ of finite combinations of the $\sech(nx)$, then $\beta(s)\ne0$ for $a<\Re s<b$.
\end{proposition}

\begin{proof}
Each approximant has Mellin transform $2\Gamma(s)\beta(s)D(s)$ with $D$ a Dirichlet polynomial. By \eqref{lv:eq:pointeval} these converge pointwise in the strip to $\Gamma(s)$, which is nonzero; a zero $\rho$ of $\beta$ in the strip would force $\Gamma(\rho)=0$.
\end{proof}

\begin{theorem}[two weight criterion]\label{lv:thm:two}
For $\frac12\le\theta<1$ the following are equivalent:
\begin{enumerate}
\item[(i)] $L(s,\chi_4)\ne0$ for $\Re s>\theta$;
\item[(ii)] $\{\sech(nx)\}_{n\in\mathbb N}$ is complete in $\mathcal W_a\cap\mathcal W_b$ for all $\theta<a<b<1$;
\item[(iii)] $S_N\to e^{-x}$ in $\mathcal W_c$ for every $c\in(\theta,1)$.
\end{enumerate}
With $\theta=\frac12$, (i) is the generalised Riemann hypothesis for $L(s,\chi_4)$.
\end{theorem}

\begin{proof}
(i) gives (a) of Proposition~\ref{lv:prop:classical}, then Theorem~\ref{lv:thm:strip} gives (iii). Convergence in each $\mathcal W_c$, $c\in[a,b]$, gives convergence in $\mathcal W_a\cap\mathcal W_b$, and Lemma~\ref{lv:lem:reduction} gives (ii). If (ii) holds and $\beta(\rho)=0$ with $\theta<\Re\rho<1$, choose $a<\Re\rho<b$ in $(\theta,1)$ and apply Proposition~\ref{lv:prop:nec}; there are no zeros with $\Re\rho\ge1$. Thus (ii) implies (i).
\end{proof}

The picture is now as follows. Completeness in one weight is free (Theorem~\ref{lv:thm:single}). Completeness in two weights straddling a strip forces that strip to be zero free (Proposition~\ref{lv:prop:nec}), and completeness for all pairs $\theta<a<b<1$ is equivalent to a zero free half plane (Theorem~\ref{lv:thm:two}); whether a zero free strip $a<\Re s<b$ alone gives completeness in $\mathcal W_a\cap\mathcal W_b$ is not shown. At a single endpoint $c<1$, convergence of $S_N$ in $\mathcal W_c$ does imply zero freeness strictly to the right of $c$, by combining it with the convergence at any $b>1$ and applying Proposition~\ref{lv:prop:nec} in $(c,b)$, but a bare zero free half plane $\Re s>c$ does not by this proof give convergence on the line $c$ itself. The properties of $\mu$ that matter form a ladder.
\begin{enumerate}
\item[(P0)] Nothing: completeness in one weight, every $c>0$.
\item[(P1)] Inversion $\sum_{d\mid N}\mu(d)=[N=1]$: the pointwise identity of Proposition~\ref{rz:prop:lvj}.
\item[(P2)] $|\mu|\le1$: the $L^1$ theorem and convergence in $\mathcal W_c$, $c>1$.
\item[(P3)] The prime number theorem for $\mu\chi_4$: convergence at $c=1$.
\item[(P4)] Power cancellation $M_\chi(x)=O(x^{\theta+\varepsilon})$: convergence and two weight completeness for $c\in(\theta,1)$, equivalent to a zero free half plane for $\beta$. Square root cancellation is GRH for $\beta$.
\end{enumerate}
The arithmetic function is the twisted $\mu\chi_4$, not $\mu$: the alternating signs of the secant produce $\beta$, not $\zeta$. The barrier $c=\frac12$ is genuine: $\beta$ vanishes at $\frac12+6.0209489\ldots i$ (computed, $6.0209489046976$, \texttt{data/lv\_sech\_tables.json}), so two weights straddling $\frac12$ can never work, and $M_\chi(x)\ne o(x^{1/2})$, since $M_\chi(x)=o(x^{1/2})$ would give $(\sigma-\frac12)/|\beta(\sigma+i\gamma)|\to0$ as $\sigma\downarrow\frac12$ at a zero, which is impossible.

\section{The zeta version}\label{lv:sec:zeta}

To involve $\mu$ itself, use the Bose kernel with its pole removed,
\begin{equation}\label{lv:eq:phi}
\varphi(x)=\frac1{e^x-1}-\frac1x,\qquad\int_0^\infty x^{s-1}\varphi(x)\,dx=\Gamma(s)\zeta(s)\quad(0<\Re s<1).
\end{equation}
Since $\varphi(x)\to-\frac12$ at the origin and $\varphi(x)\sim-1/x$ at infinity, $\varphi\in\mathcal W_c$ exactly for $0<c<1$. The cosecant $\operatorname{csch}x$, with transform $2\Gamma(s)(1-2^{-s})\zeta(s)$, is unsuitable because the pole of $\zeta$ removes it from $\mathcal W_c$ for $c\le1$. Write $T_N(x)=\sum_{n\le N}\mu(n)\varphi(nx)$.

\begin{theorem}\label{lv:thm:zeta}
(1) For every $c\in(0,1)$ the dilates $\{\varphi(nx)\}_{n\in\mathbb N}$ are complete in $\mathcal W_c$, unconditionally. (2) Let $\frac12\le\theta<1$ and suppose $M(x)=O_\varepsilon(x^{\theta+\varepsilon})$ for every $\varepsilon>0$. Then $T_N\to e^{-x}$ in $\mathcal W_c$ with $\|T_N-e^{-x}\|_c=O_\varepsilon(N^{\theta-c+\varepsilon})$ for $\theta<c<1$ and $0<\varepsilon<c-\theta$. (3) For $\frac12\le\theta<1$, completeness of $\{\varphi(nx)\}$ in $\mathcal W_a\cap\mathcal W_b$ for all $\theta<a<b<1$ is equivalent to $\zeta(s)\ne0$ for $\Re s>\theta$ (the pole at $1$ being allowed); with $\theta=\frac12$ it is equivalent to the Riemann hypothesis.
\end{theorem}

\begin{proof}
The proofs of Theorems~\ref{lv:thm:single}, \ref{lv:thm:strip}, \ref{lv:thm:two} carry over with $(\beta,\mu\chi_4,\sech)$ replaced by $(\zeta,\mu,\varphi)$, with the following changes. The Mellin transform of $\varphi(n\,\cdot)$ is $n^{-s}\Gamma(s)\zeta(s)$ on $0<\Re s<1$, with no factor $2$, and $|\zeta(c+it)|\ll(1+|t|)^{1/2}$ there, so $|\Gamma\zeta|\le C(1+|t|)^Ae^{-\pi|t|/2}$ and the proof of part (1) is that of Theorem~\ref{lv:thm:single}; $\zeta$ has discrete zeros on the line. For part (2) the Mellin transform of $T_N-e^{-x}$ on the line $\Re s=c$ is $\Gamma(s)\bigl(\zeta(s)\mathcal A_N(s)-1\bigr)$ with $\mathcal A_N=\sum_{n\le N}\mu(n)n^{-s}$, and Lemma~\ref{lv:lem:ps} holds for $\mathcal A_N$ and $1/\zeta$ with $M$ in place of $M_\chi$ ($1/\zeta$ is analytic on $\Re s>\theta$ because $\zeta$ has no zeros there, and $\zeta$ is analytic on $0<\Re s<1$). For part (3) Proposition~\ref{lv:prop:nec} holds with $\zeta$ because the approximants have Mellin transform $\Gamma\zeta D$ and converge pointwise in the strip $a<\Re s<b\subset(0,1)$ to $\Gamma$; zeros with $\Re\rho\ge1$ do not exist; and the classical equivalence of Proposition~\ref{lv:prop:classical} for $\zeta$ is Littlewood's theorem.
\end{proof}

The contrast with the secant is that $\varphi$ has no unconditional M\"obius range at all: its Mellin line must lie inside $0<c<1$. Part (3) is a Nyman and Beurling type criterion on the half line, and part (1) explains why it needs two weights: in the classical Nyman and Beurling problem the depth comes from restricting to $(0,1)$, which makes Mellin transforms analytic in a half plane, while here the second weight plays that role.

\section{What the quasi Riemann hypothesis gives}\label{lv:sec:qrh}

In October 2026 OpenAI released a preprint \cite{OAI26a} (dated 30 September) announcing a zero free half plane $\Re s>\frac78$ for Dirichlet $L$ functions, with a formalisation in Lean, and a companion \cite{OAI26b} with the weaker half plane $\Re s>\frac{11}{12}$; Liu \cite{Liu26a} claims a small improvement of the exponent $\frac78$. A preprint attributed to Thorner, extending a zero free half plane $\Re s>\frac{31}{32}$ to the Dedekind zeta function of every number field, is cited here from the source as unchecked \cite{Thorner26}. None of these has been refereed, and none has been checked for this book. Every statement below that uses them is conditional on them, and those about $\beta$ on one further point: the $\frac78$ theorem is announced for Dirichlet $L$ functions, and that it covers $L(s,\chi_4)$ is assumed here and is a question for the authors. If these results stand, the functional equation puts the nontrivial zeros of $\zeta$ and of $\beta$ in $\frac18\le\Re\rho\le\frac78$.

\begin{center}\footnotesize
\begin{tabular}{p{0.32\textwidth}p{0.3\textwidth}p{0.28\textwidth}}
\toprule
input & secant, $\mu\chi_4$ & zeta kernel, $\mu$\\
\midrule
none & one weight, all $c>0$ & one weight, $0<c<1$\\
PNT & M\"obius series, $c\ge1$ (rate subexponential in $\sqrt{\log N}$) & nothing new\\
Dedekind QRH ($\frac{31}{32}$), unchecked & $c>\frac{31}{32}$ & $\frac{31}{32}<c<1$\\
QRH for Dirichlet $L$ ($\frac78$), conditional & $\frac78<c$ & $\frac78<c<1$\\
GRH & $c>\frac12$ & $\frac12<c<1$\\
\bottomrule
\end{tabular}
\end{center}

\noindent For a zero free half plane $\Re s>\theta$ the entries mean: the M\"obius series converges in $\mathcal W_c$ at the rate $N^{\theta-c+\varepsilon}$ for every $c>\theta$ (secant) or $\theta<c<1$ (zeta kernel), and two weight completeness holds on that range. For the secant kernel and $c>1$ the elementary absolute convergence argument gives $O_c(N^{1-c})$, and at $c=1$ the prime number theorem gives $O((1+\sqrt{\log N})e^{-a\sqrt{\log N}})$; neither is a power law with $\theta=1$, so the PNT row is not an instance of the rate $N^{\theta-c+\varepsilon}$. The input $\frac78$ pushes the M\"obius route below $c=1$, from $c\ge1$ down to $c>\frac78$, for both kernels at once if the theorem covers $\beta$ as well as $\zeta$. The stated half plane does not supply convergence in the remaining range $\frac12<c\le\frac78$; GRH is sufficient there, whereas a weaker zero free bound already gives convergence above its own boundary. For the one weight question QRH is irrelevant, since that question was settled unconditionally.

\section{Relation to the Thorin programme}\label{lv:sec:thorin}

Three remarks connect this part to the rest of the book. The second and third are analogies (heuristic), not consequences of the theorems above.

First, the secant is a Thorin object, in the form of Lemma~\ref{rz:lem:sech}: $\sech\sqrt{2s}=\prod_{k\ge1}(1+s/\lambda_k)^{-1}$ with $\lambda_k=(2k-1)^2\pi^2/8$ is the Laplace transform of the exit time of Brownian motion from $(-1,1)$, a generalised gamma convolution with discrete Thorin measure at the $\lambda_k$ (Table~\ref{tab:ch3:cat}). Dilation of $\sech$ corresponds to a time scaling of that law after the substitution $x=\sqrt{2s}$.

Second, the book's Thorin type expansion of $1/\zeta(\frac12+\sqrt s)$ (Theorem~\ref{rz:thm:sech}) applies Proposition~\ref{rz:prop:lvj} pointwise at $x=\sqrt s\log n$ to the Dirichlet series of $1/\zeta$. Theorems~\ref{lv:thm:two} and~\ref{lv:thm:zeta} are statements in $L^2$ in the variable $x$ about approximating $e^{-x}$ in weighted spaces; they are not statements about that series in $s$, and the kernel $\varphi$ is not used in it. The parallel is this: the convergence of the construction needs the same zero free region as the M\"obius series, and no new zero free region comes out of it.

Third, the dichotomy between one and two weights is a pattern for the programme. Statements made on a single Mellin line (completeness, $L^2$ approximation) are cheap, because point evaluations off the line are not controlled, while statements uniform on a strip carry analyticity into the strip and force zero freeness there (Proposition~\ref{lv:prop:nec}). A GGC Laplace transform is analytic and zero free in a half plane, so a Thorin representation of $1/\xi$ is of the second kind; Section~\ref{lv:sec:thorinform} makes this precise for $\xi(s,\chi_4)$.

\section{Numerical check}\label{lv:sec:numcheck}

The weighted errors $\|S_N-e^{-x}\|_c$ and $\|T_N-e^{-x}\|_c$ were computed by quadrature on logarithmic grids (\texttt{code/lv\_\allowbreak sech\_\allowbreak tables.py}, output \texttt{data/lv\_\allowbreak sech\_\allowbreak tables.json}); the secant table was computed on two grid steps and agrees to $10^{-15}$ relative, the zeta kernel table to $5\cdot10^{-8}$, the part of the zeta integrals beyond $x=60$ being added exactly ($T_N-e^{-x}=-m_N/x$ there, $m_N=\sum_{n\le N}\mu(n)/n$). The pointwise identity of Proposition~\ref{rz:prop:lvj} is proved, and the check of it at $x=0.05,0.3,1,3$ in the source is not reproduced here.

\begin{table}[h]\centering\footnotesize
\begin{tabular}{rcccc}
\toprule
& \multicolumn{4}{c}{secant, $\|S_N-e^{-x}\|_c$}\\
$N$ & $c=0.60$ & $c=0.75$ & $c=1.00$ & $c=1.25$\\
\midrule
$10$ & $7.3\times10^{-3}$ & $5.1\times10^{-3}$ & $2.9\times10^{-3}$ & $1.7\times10^{-3}$\\
$10^2$ & $3.9\times10^{-2}$ & $1.9\times10^{-2}$ & $6.3\times10^{-3}$ & $2.2\times10^{-3}$\\
$10^3$ & $3.5\times10^{-2}$ & $1.2\times10^{-2}$ & $2.2\times10^{-3}$ & $4.2\times10^{-4}$\\
$10^4$ & $5.3\times10^{-2}$ & $1.3\times10^{-2}$ & $1.2\times10^{-3}$ & $1.3\times10^{-4}$\\
\midrule
& \multicolumn{3}{c}{zeta kernel, $\|T_N-e^{-x}\|_c$}\\
$N$ & $c=0.75$ & $c=0.875$ & $c=0.95$\\
\midrule
$10$ & $2.1\times10^{-1}$ & $2.3\times10^{-1}$ & $3.1\times10^{-1}$\\
$10^2$ & $1.2\times10^{-1}$ & $1.0\times10^{-1}$ & $1.2\times10^{-1}$\\
$10^3$ & $3.0\times10^{-2}$ & $1.9\times10^{-2}$ & $1.9\times10^{-2}$\\
$10^4$ & $2.6\times10^{-2}$ & $1.2\times10^{-2}$ & $1.0\times10^{-2}$\\
\bottomrule
\end{tabular}
\caption{Weighted errors of the M\"obius approximants (computed). The column $c=0.95$ of the source was $2.9\times10^{-1}$, $1.1\times10^{-1}$, $1.8\times10^{-2}$, $9.8\times10^{-3}$, a few percent low.}\label{lv:tab:secant}
\end{table}

For $c\ge1$ Theorem~\ref{lv:thm:c1} gives convergence but no rate, and at $c=1$ the computed errors rise from $N=10$ to $N=100$. Inside the strip the GRH rates $N^{1/2-c}$ are slow ($N^{-0.25}$ at $c=0.75$, $N^{-0.1}$ at $c=0.6$) and are masked at this range by the oscillation of $M_\chi(N)$, whose values at $N=10,10^2,10^3,10^4$ are $2,3,6,-24$; the error is concentrated at $x\lesssim1/N$, where $S_N(x)\approx\frac12M_\chi(N)$ (heuristic, a Taylor approximation). The finite computations cannot prove an asymptotic rate or a zero free half plane. The largest value of $|M_\chi(n)|/\sqrt n$ for $n\le10^4$ is $2/\sqrt3=1.1547$, attained at $n=3$, so it says nothing about large $n$. The zeta kernel series converges only conditionally pointwise, since $\varphi(nx)\approx-1/(nx)$ and $\sum\mu(n)/n=0$ is the prime number theorem; its weighted errors nevertheless decrease steadily.
